\chapter{NÚMEROS ÍNDICES}

\section{Introdução}

O conceito de números-índices pode ser aplicado a qualquer
variável ou conjunto de variáveis que se deseje comparar ao
longo do tempo ou ao longo do espaço (comparações regionais). São instrumentos fundamentais da estatística aplicados à análise comparativa de fenômenos cuja função é transformar valores absolutos em medidas relativas, permitindo a visualização de variações de forma padronizada.\vskip0.3cm

Tradicionalmente, o índice só tem significado se as datas a que se refere forem claramente especificadas, caracterizando a época, o período e o local a que dizem respeito.\vskip0.3cm


No contexto da gestão pública, especialmente na área de trânsito, os números-índices são valiosos para monitorar a eficiência de serviços, a evolução de indicadores de mobilidade, a incidência de acidentes e o comportamento dos condutores, entre outros aspectos relevantes para o planejamento urbano.


\section{Índice Simples (ou Individual)}

O Índice Simples mede a variação relativa de uma única variável em relação a um período base.

\[
I(P_{t}/P_{0}) = \frac{P_{t}}{P_{0}} \times 100
\]

Onde:
\begin{itemize}
  \item \( $P_{t}$ \) é o valor no período atual;
  \item \( $P_{0}$\) é o valor no período base.
\end{itemize}

Se \( I_t = 120 \), houve um aumento de 20\% em relação ao período base.



\section{Índice Relativos}



\section{Índice Compostos}

\subsection{Indice de Laspeyres}
\subsection{Indice de Paasche}






















